\documentclass[10pt,a4paper]{amsart}
\usepackage[margin=19mm]{geometry}
\usepackage{amsmath,amssymb,mathtools}
\usepackage[scr=rsfs,cal=euler]{mathalfa}
\usepackage{xfrac}
\usepackage[unicode,hidelinks,bookmarks=false]{hyperref}
\newcommand{\ie}{i.e.\ }
\newcommand{\cf}{cf.\ }
\newcommand{\resp}{respectively\ }
\newcommand{\Subsection}{\subsection}
\providecommand{\nobreakdash}{\nobreak\hbox{-}}
\setlength{\emergencystretch}{2em}
\multlinegap0pt
\usepackage{bm}
\usepackage{bbm}
\usepackage{dsfont}
\usepackage{xspace}
\usepackage{cancel}
\usepackage{mathtools}
\usepackage{amsthm}
\makeatletter
\@ifundefined{proposition}{\newtheorem{proposition}{Proposition}}{}
\makeatother
\DeclareMathOperator{\eps}{\varepsilon}
\title[A symmetric mechanism for symmetry-breaking in oscillator ne]{A symmetric mechanism for symmetry-breaking in oscillator networks with strong nonlinear coupling}
\author{Theodore Vo, Yangyang Wang}
\date{Source: arXiv:2606.17780v1 (CC BY 4.0)}
\begin{document}
\maketitle

\noindent\emph{Source and licence.} A symmetric mechanism for symmetry-breaking in oscillator networks with strong nonlinear coupling. Version arXiv:2606.17780v1 (CC BY 4.0), distributed under the Creative Commons Attribution 4.0 International licence, as recorded in the arXiv licence metadata for this exact version and shown on the abstract page https://arxiv.org/abs/2606.17780v1. Full rights notice: licenses/arxiv-2606.17780-CC-BY-4.0.txt.\par\smallskip

\noindent\textit{Editorial scope.} Counted unit: the Proposition whose source label is \texttt{eq:dynamicsonWc} in the article (source0.tex lines 574--599, source label \texttt{eq:dynamicsonWc}), delivered together with the article's own complete original proof of it (source0.tex lines 601--652, one \texttt{proof} environment of 52 lines). No mathematics is written, repaired or re-derived: statement and proof are the article's own text line for line. Required context retained uncounted: eq:normalform (lines 201--208); eq:nfsymasym (lines 550--558). Every definition, display and label of those lines is retained unchanged. Omitted from this package and not redistributed: 1--200, 209--549, 559--573, 600--600, 653--2238. Nothing inside a retained excerpt is omitted or altered: every retained body is delivered exactly as the article's lines are written, so no omission receipt of any kind accompanies this package. Citations: the retained excerpts carry no citation command at all, so no bibliography entry of the article is transcribed and no reference is redistributed. Reference adaptation: none was needed and none was made; every reference inside the delivered unit resolves inside it, so no unresolved reference or citation is produced. Printed slips kept literal: no printed word, symbol, hypothesis or number of the counted statement or of its proof was altered. Layout and class: the article's own class is \texttt{article} with packages including geometry, graphicx, caption, hyperref, authblk, enumitem, grffile, parskip, amssymb, amsmath, amsthm, xcolor, mathrsfs, footmisc; the wrapper is the lane's pinned \texttt{amsart} editorial wrapper at 10pt on A4, which restates the theorem-like environments the retained text needs and adds the article's own macro definitions that the retained text needs (\texttt{\textbackslash eps}), with extra wrapper packages (bm, bbm, dsfont, xspace, cancel, mathtools). The delivered numbering is the unit's own. Assets: no retained span contains a figure environment, a graphics-inclusion command, a PDF-page inclusion, a table or a TikZ picture; no image file is copied into this package, the package declares no asset dependency and no image file is redistributed with it. Licence: version arXiv:2606.17780v1 of this article is distributed under the Creative Commons Attribution 4.0 International licence, as recorded in the arXiv licence metadata for this exact version and shown on the abstract page https://arxiv.org/abs/2606.17780v1. A full rights notice is delivered at licenses/arxiv-2606.17780-CC-BY-4.0.txt. Characters: every retained excerpt is pure ASCII, so the delivered engine, XeLaTeX with its default Latin Modern text font, reports no missing character.
\begin{equation} \label{eq:normalform}
  \begin{split}
    \dot u_1 &= -u_1-u_2+v_1+u_1^2-\left( 1+\tfrac{1}{2}\lambda_1 \right) u_1 v_2 -a (u_1 v_2+u_2 v_1) + P_1(u_1,u_2,v_1,v_2,\eps), \\
    \dot u_2 &= -u_2-u_1+v_2+u_2^2  - \left( 1+\tfrac{1}{2}\lambda_1 \right) u_2 v_1 -a (u_2 v_1+u_1 v_2)+ P_2(u_2,u_1,v_2,v_1,\eps), \\
    \dot v_1 &= \eps \left[ 1+\tfrac{1}{2} (\lambda_1+\lambda_2) u_1+ b u_1^2 + \tfrac{1}{4} a (\lambda_1+\lambda_2)u_2^2 + Q_1(u_1,u_2,v_1,v_2,\eps) \right], \\ 
    \dot v_2 &= \eps \left[ 1+\tfrac{1}{2} (\lambda_1+\lambda_2) u_2 + b u_2^2+ \tfrac{1}{4} a \left( \lambda_1+\lambda_2 \right) u_1^2  + Q_2(u_2,u_1,v_2,v_1,\eps) \right],
  \end{split}
\end{equation}

\begin{equation} \label{eq:nfsymasym}
  \begin{split}
    \frac{d\hat u_1}{dt} &= -2\hat u_1 + \hat v_1 + \hat u_1^2+\hat u_2^2 - \tfrac{4a+2+\lambda_1}{2} \hat u_1 \hat v_1 + \tfrac{2+\lambda_1}{2}\hat u_1 \hat v_2+\tfrac{4a+2+\lambda_1}{2} \hat u_2 \hat v_2 + k \eps \hat u_1 + \hat P_1, \\
    \frac{d\hat u_2}{dt} &= \hat v_2+2 \hat u_1 \hat u_2-\tfrac{2+\lambda_1}{2} \hat u_2 \hat v_1+k \eps \hat u_2 + \hat P_2, \\
    \frac{d\hat v_1}{dt} &= \eps \left[ 1+\tfrac{\lambda_1+\lambda_2}{2} \hat u_1 + \tfrac{4b+a(\lambda_1+\lambda_2)}{4} \hat u_1^2 + \tfrac{4b-a(\lambda_1+\lambda_2)}{2} \hat u_1 \hat u_2 + \tfrac{4b+a(\lambda_1+\lambda_2)}{4} \hat u_2^2 + \hat Q_1 \right], \\
    \frac{d\hat v_2}{dt} &= \eps \left[ \tfrac{\lambda_1+\lambda_2}{2} \hat u_2 +\tfrac{4b-a(\lambda_1+\lambda_2)}{2} \hat u_1 \hat u_2 + \hat Q_2 \right], \\
    \frac{d\eps}{dt} &= 0,
  \end{split}
\end{equation}

\begin{proposition}
For the system \eqref{eq:nfsymasym}, there exists an attracting 4D center manifold of the origin given by 
%
\[ W^c = \left\{ (u_1,v_1,u_2,v_2,\eps) : 
u_1 = h(u_2,v_1,v_2,\eps)
\right\}, \]
%
where the function $h$ is given by 
%
\[ h = \tfrac{1}{2}u_2^2 -\tfrac{4a+1+\lambda_1}{2}v_1^2 - \tfrac{4a+\lambda_1}{8} v_2^2+\tfrac{4a+\lambda_1}{4} u_2 v_2 + \tfrac{2+\lambda_1}{4} v_1 v_2 +\tfrac{4a+4k-\lambda_2}{4} \eps v_1 - \tfrac{2+\lambda_1}{8} \eps v_2 + \mathcal O \left( \eps^2,(u_2+v_1+v_2)^3 \right). \]
%
Moreover, the reduced dynamics on the center manifold are given by
%
\begin{equation} \label{eq:dynamicsonWc}
  \begin{split}
    %\dot u_2 &= v_2 - \lambda_1 u_2 v_1 + \tfrac{4k-2-\lambda_1}{2} \eps u_2 + \tfrac{2+5a+2a^2-b}{2}u_2^3 + \mathcal O \left(u_2^2 v_2, u_2 v_1^2, u_2 v_2^2 \right) + \eps \mathcal O \left( u_2 v_1, u_2 v_2, v_1 v_2 \right), \\
    \dot u_2 &= v_2 - \lambda_1 u_2 v_1 + \gamma_1 u_2 v_1^2 + \gamma_2 u_2^3 + \gamma_3 \eps z + \mathcal O \left( u_2^2 v_2, u_2 v_2^2  \right) + \eps \mathcal O \left( u_2 v_1, u_2 v_2, v_1 v_2 \right), \\
    %\dot v_1 &= \eps \left[ 1+\tfrac{1}{2}(\lambda_1+\lambda_2)v_1 + \tfrac{4\mu+(1+a)(\lambda_1+\lambda_2)}{4} u_2^2 + \mathcal O \left(u_2v_1,u_2v_2,v_1^2, v_2^2, \eps^2 \right)  \right], \\ 
    \dot v_1 &= \eps \left[ 1 + \tfrac{\lambda_1+\lambda_2}{2} v_1 + \alpha_1 v_1^2 + \alpha_2 u_2 v_1 + \alpha_3 u_2^2 + \mathcal O \left( u_2 v_2, v_2^2, \eps^2 \right) \right], \\
    %\dot v_2 &= \eps \left[ (\lambda_1+\lambda_2) u_2 + \mathcal O\left( u_2 v_1, v_1v_2,\eps^2 \right) \right].
    \dot v_2 &= \eps \left[ (\lambda_1+\lambda_2) u_2 + \beta u_2 v_1 + \mathcal O \left( v_1 v_2, \eps^2 \right) \right],
  \end{split}
\end{equation}
%
where the coefficients $(\alpha_1,\alpha_2,\alpha_3,\beta,\gamma_1,\gamma_2,\gamma_3)$ can be computed in terms of the coefficients of \eqref{eq:normalform}.
\end{proposition}

\begin{proof}
The fixed point at the origin of the extended system \eqref{eq:nfsymasym} has stable spectrum $\sigma_s = \{ -2 \}$ and center spectrum $\sigma_c = \{ 0,0,0,0 \}$, with corresponding stable and center subspaces, respectively, given by
%
\[ \mathbb E^s(0) = \operatorname{span} \begin{bmatrix} 1 \\ 0 \\ 0 \\ 0 \\ 0 \end{bmatrix} \quad {\rm and } \quad \mathbb E^c(0) = \operatorname{span}  \left\{ 
\begin{bmatrix} 1 \\ 0 \\ 2 \\ 0 \\ 0 \end{bmatrix} , 
\begin{bmatrix} 0 \\ 1 \\ 0 \\ 0 \\ 0 \end{bmatrix} , 
\begin{bmatrix} 0 \\ 0 \\ 0 \\ 1 \\ 0 \end{bmatrix} ,
\begin{bmatrix} 0 \\ 0 \\ \tfrac{1}{2} \\ 0 \\ 1 \end{bmatrix} 
\right\}. \]
%
We align the center directions with the coordinate axes via the change of coordinates
%
\[ \widetilde u_1 = \hat u_1-\tfrac{1}{2} \hat v_1 + \tfrac{1}{4} \eps, \quad \widetilde v_1 = \tfrac{1}{2}\hat v_1 - \tfrac{1}{4} \eps, \quad \widetilde u_2 = \hat u_2, \quad \widetilde v_2 = \hat v_2, \quad \text{ and } \quad \widetilde \eps = \tfrac{1}{2} \eps. \]
%
Transforming and dropping tildes, the linear part of the vector field has coefficient matrix
%
%\begin{equation}  \label{eq:nfjordan}
%  \begin{split}
%    \begin{bmatrix} \dot u_1 \\ \dot v_1 \\ \dot u_2 \\ \dot v_2 \\ \dot \eps \end{bmatrix} &= 
%    \begin{bmatrix} -2 & 0 & 0 & 0 & 0 \\ 0 & 0 & 0 & 0 & 1 \\ 0 & 0 & 0 & 1 & 0 \\ 0 & 0 & 0 & 0 & 0 \\ 0 & 0 & 0 & 0 & 0 \end{bmatrix}
%    \begin{bmatrix} u_1 \\ v_1 \\ u_2 \\ v_2 \\ \eps \end{bmatrix} +
%    \begin{bmatrix} u_1^2+u_2^2-\lambda_1 u_1 v_1 - (1+\lambda_1) v_1^2+(1+\tfrac{1}{2}\lambda_1) u_2 v_2 + H_{1}  \\ \eps \left( \tfrac{1}{2}(\lambda_1+\lambda_2) (u_1+v_1) + K_{1} \right) \\ 2 u_1 u_2 + (1+\tfrac{1}{2}\lambda_1) u_1 v_2 - \lambda_1 u_2 v_1 + (1+\tfrac{1}{2}\lambda_1) v_1 v_2 + H_{2} \\ \eps \left( (\lambda_1+\lambda_2)u_2 + K_{2} \right) \\ 0 \end{bmatrix},
%  \end{split}
%\end{equation}
\[ Df = \begin{bmatrix} -2 & 0 & 0 & 0 & 0 \\ 0 & 0 & 0 & 0 & 1 \\ 0 & 0 & 0 & 1 & 0 \\ 0 & 0 & 0 & 0 & 0 \\ 0 & 0 & 0 & 0 & 0 \end{bmatrix}. \]
%
%where $H_{1}, H_{2}, K_{1}$, and $K_{2}$ denote the higher order terms. 
By the Center Manifold Theorem, there exists a 4D center manifold, $W^c$, of the origin given by 
%
\[ W^c = \left\{ (u_1,u_2,v_1,v_2,\eps) : u_1 = h(u_2,v_1,v_2,\eps) \right\} \]
%
where the function $h$ satisfies the invariance equation 
%
\begin{equation*} 
  \begin{split}
\tfrac{\partial h}{\partial u_2} \dot u_2 + \tfrac{\partial h}{\partial v_1} \dot v_1 + \tfrac{\partial h}{\partial v_2} \dot v_2 = -2h &+ h^2+u_2^2-(1+\lambda_1+4a)v_1^2-(4a+\lambda_1) h v_1  \\ 
& + \tfrac{2+\lambda_1}{2} v_2 \left( h+\tfrac{4a+2+\lambda_1}{2+\lambda_1}u_2 + v_1 \right) + \tfrac{4k-(2+\lambda_1+4a+\lambda_1+\lambda_2)}{2} \eps (h+v_1) + H,
  \end{split}
\end{equation*}
%
and $H$ denotes the higher-order terms.
By expanding $h$ as a power series of the form $\displaystyle h = \sum_k h_k$, where $h_k$ is a homogeneous polynomial of degree $k$ in the variables $(u_2,v_1,v_2,\eps)$, and equating coefficients of like terms, we obtain the function $h$ given in the statement of the proposition. 
%
%\begin{equation*}
%  \begin{split}
%   h &= \tfrac{1}{2}u_2^2-\tfrac{1}{2}(1+\lambda_1)v_1^2-\tfrac{1}{8}\lambda_1 u_2v_2-\tfrac{1}{4} \lambda_2 \eps v_1 + \tfrac{1}{8} \lambda_2 \eps^2 + H,
%  \end{split}
%\end{equation*}
%
%where $H$ denotes terms of cubic or higher order. 
The reduced dynamics on the center manifold are then obtained by substituting $u_1 = h$ into the ODEs and relabelling the coefficients.
\end{proof}

\end{document}
